\documentclass[12pt]{article}
\usepackage{../style_files/aas_macros}
\usepackage{natbib}
\usepackage[utf8]{inputenc}
\usepackage[left=3cm,right=3cm,top=3cm,bottom=3cm]{geometry}
\usepackage{color}
\usepackage[colorlinks = true,
            linkcolor = blue,
            urlcolor  = blue,
            citecolor = blue]{hyperref}
\usepackage{graphicx}
\usepackage{subfigure}
\usepackage{tablefootnote}
\usepackage{longtable}
\usepackage{tabularx}
\usepackage{booktabs}
\usepackage{makecell}
\usepackage{soul}
\sethlcolor{lightgray}
\usepackage[most]{tcolorbox}

\tcbset{textmarker/.style={%
        enhanced,
        parbox=false,boxrule=0mm,boxsep=0mm,arc=0mm,
        outer arc=0mm,left=6mm,right=3mm,top=7pt,bottom=7pt,
        toptitle=1mm,bottomtitle=1mm,oversize}}

\newtcolorbox{hintBox}{textmarker,
    borderline west={6pt}{0pt}{yellow},
    colback=yellow!10!white}
\newcommand{\todo}[1]{\begin{hintBox} \textbf{TODO:} #1 \end{hintBox}}

\hypersetup{citecolor=blue}

\definecolor{MSBlue}{rgb}{.204,.353,.541}
\usepackage{titlesec}
\titleformat*{\section}{\color{MSBlue}}
\titleformat*{\subsection}{\color{MSBlue}}
\setlength{\parindent}{0pt}

\title{Gamma/Hadron Separation}
\author{Gernot Maier}
\date{\today}

\begin{document}

\maketitle

For gamma/hadron separation, energy-binned XGBoost classifiers are trained on
Monte Carlo simulations of gamma-ray and proton-initiated air showers.

For gamma/hadron separation, energy-binned XGBoost classifiers replace the
TMVA-based BDTs used in the classical analysis.

\section{Gamma/Hadron Separation}
\label{sec:classification}

\subsection{Model architecture}

Gamma/hadron separation is implemented as a binary XGBoost classifier
(\texttt{XGBClassifier}, objective \texttt{binary:logistic})
that assigns each event a probability $p_\gamma \in [0,1]$ of being
a gamma-ray-initiated shower.
Because the classification problem is energy-dependent (shower morphology and
background composition change strongly with energy), separate models are trained
in nine energy bins spanning $\log_{10}(E_\mathrm{rec}/\mathrm{TeV}) \in [-2, 2.5]$.
During application, adjacent bin models are interpolated linearly in
$\log_{10}(E_\mathrm{rec})$ using interpolation weights $\alpha$:
\begin{equation}
  p_\gamma = (1 - \alpha)\, p_{\gamma,\mathrm{lo}} + \alpha\, p_{\gamma,\mathrm{hi}}.
\end{equation}

The classification features are a subset of the reconstruction features
(Hillas parameters, reduced stereo quantities, MSCW/MSCL shape parameters)
listed in Appendix~\ref{sec:appendix_features}.
The energy estimate is used only for binning and is not a classifier input.

Hyperparameters for the classifier are listed in Table~\ref{tab:hyperparameters_cls}.
A lower tree depth (4) compared to the regression model avoids overfitting given the
binary label noise inherent in hadronic background samples.

\begin{table}[h]
\centering
\caption{XGBoost classification hyperparameters.}
\label{tab:hyperparameters_cls}
\begin{tabular}{ll}
\toprule
\textbf{Parameter} & \textbf{Value} \\
\midrule
\texttt{objective}           & \texttt{binary:logistic} \\
\texttt{eval\_metric}        & log-loss, AUC \\
\texttt{n\_estimators}       & 5\,000 (with early stopping) \\
\texttt{early\_stopping\_rounds} & 100 \\
\texttt{learning\_rate}      & 0.02 \\
\texttt{max\_depth}          & 4 \\
\texttt{gamma}               & 0.2 \\
\texttt{subsample}           & 0.8 \\
\texttt{colsample\_bytree}   & 0.6 \\
\bottomrule
\end{tabular}
\end{table}

\subsection{Signal efficiency and hadron rejection}

The classifier output $p_\gamma$ is converted to a binary label by applying a
threshold $t_\gamma$ calibrated to a target signal efficiency
(e.g., $\epsilon_\gamma = 80\%$).
The threshold is determined from the signal and background efficiency curves
computed on the test sample.
Thresholds at a range of signal efficiencies (20\%--95\% in 5\% steps) are
stored with the trained model for flexible analysis cuts.
The classification performance is also characterised by the ROC curve
(gamma efficiency vs.\ hadron rejection $1 - \epsilon_h$),
the Q-factor $Q = \epsilon_\gamma / \sqrt{\epsilon_h}$,
and the distribution of the classifier score for signal and background.

\subsection{Diagnostic checks for classification}

\todo{The following plots are generated by \texttt{plot\_classification\_performance\_metrics.py}. Produce one figure per energy bin (9 bins) plus one combined summary figure.}

\paragraph{Signal and background efficiency curves.}
Signal efficiency $\epsilon_\gamma(t)$ and background efficiency $\epsilon_h(t)$
as a function of the classifier threshold $t$,
shown for each energy bin and for the XGBoost and TMVA-BDT classifiers.
\todo{Figure: $\epsilon_\gamma$ and $\epsilon_h$ vs.\ threshold, XGBoost (dashed) vs.\ TMVA-BDT (solid), one panel per energy bin.}

\paragraph{ROC curve.}
$\epsilon_\gamma$ vs.\ $1 - \epsilon_h$ for each energy bin,
with the AUC value quoted as a summary metric.
\todo{Figure: ROC curves per energy bin, XGBoost vs.\ TMVA-BDT; AUC values in legend.}

\paragraph{Q-factor curve.}
$Q(\epsilon_\gamma) = \epsilon_\gamma / \sqrt{\epsilon_h}$ as a function of gamma efficiency,
for each energy bin.
The maximum Q-factor provides a single-number comparison between classifiers.
\todo{Figure: Q-factor vs.\ $\epsilon_\gamma$ per energy bin, with maximum Q quoted for XGBoost and TMVA-BDT.}

\paragraph{Score distributions.}
Probability density of $p_\gamma$ for signal (gamma MC) and background (proton MC),
reconstructed from the derivative of the efficiency curves.
Overlap between the distributions indicates classification difficulty.
\todo{Figure: Score distributions (signal and background PDF) per energy bin, XGBoost vs.\ TMVA-BDT.}

\paragraph{SHAP feature importance for classification.}
Top-20 features by mean absolute SHAP value for the \texttt{label} target,
shown for a representative energy bin.
\todo{Figure: SHAP importance bar chart for classification (one representative bin plus summary across all bins). Generated by \texttt{diagnostic\_shap\_summary.py}.}

For gamma/hadron separation, energy-binned XGBoost classifiers with linear
interpolation between adjacent bins replace the classical TMVA-BDT analysis.

% TODO: Summarise quantitative improvements in angular resolution, energy resolution,
% energy bias, and gamma/hadron separation quality factor once result plots are available.
\todo{Fill in quantitative summary once result plots are generated: angular resolution improvement (\%), energy resolution improvement (\%), energy bias reduction, max Q-factor comparison (XGBoost vs.\ TMVA-BDT), and Li\&Ma significance gain on Crab data.}

\section{Detailed description of features}
\label{sec:appendix_features}

\subsection{Additional features for gamma/hadron classification}

\begin{longtable}{lll p{5.5cm}}
\caption{Features used exclusively for gamma/hadron classification.} \label{tab:features-class} \\
\hline
\textbf{Feature} & \textbf{Unit} & \textbf{Range (clipped)} & \textbf{Description} \\
\hline
\endfirsthead
\hline
\textbf{Feature} & \textbf{Unit} & \textbf{Range (clipped)} & \textbf{Description} \\
\hline
\endhead
\hline
\endfoot
\texttt{MSCW} & -- & $[-2,2]$ & Mean scaled image width \\
\texttt{MSCL} & -- & $[-2,5]$ & Mean scaled image length \\
\texttt{EmissionHeightChi2} & log$_{10}$ & $[-6,4]$ & $\chi^2$ of emission height fit \\
\texttt{EChi2S} & log$_{10}$ & $[-6,4]$ & $\chi^2$ of stereo energy fit \\
\texttt{SizeSecondMax} & log$_{10}$ p.e. & $[0,5]$ & Size of the second-largest image \\
\texttt{DispAbsSumWeigth} & -- & $[0,5]$ & Sum of DispBDT weights \\
\end{longtable}


\section*{Acknowledgments}

Add acknowledgments here.

%%%%%%%%%%%%%%%%%%%%%%%%%%%%%%%%%%%%%%%%%%%
\bibliographystyle{aasjournal}
\bibliography{../bibliography/bibliography.bib}


\end{document}
